%% ma: L12-B14-0121
%% lop: 12
%% chuong: 5
%% bai: 14
%% dang: 12.14.32
%% loai: TLN
%% muc: VD
%% dap_an: "2,5"
%% nguon: "Ảnh giáo viên gửi 07/11/2026 (Tư Duy Mở - tuduymo.com, đề tự luyện Phương trình mặt phẳng, đề khoảng cách), Câu 17"
%% kien_thuc: "Bài 14 (hai mặt phẳng song song, khoảng cách từ một điểm đến mặt phẳng)"
%% trang_thai: cho-duyet
%% ly_do: "Dạng 12.14.32 mới đề xuất, chờ giáo viên duyệt."
%% ghi_chu: "Tài liệu gốc không kèm đáp án; đáp án do Claude giải và kiểm bằng sympy. Đề gốc không ghi điều kiện (Q) song song với (P) nhưng có dùng khái niệm khoảng cách giữa hai mặt phẳng (P) và (Q), chỉ xác định khi hai mặt phẳng song song; đã ghi rõ điều kiện."
\begin{ex}
Trong không gian $Oxyz$, cho mặt phẳng $(P)\colon x-2y+2z-11=0$ và điểm $A(1;0;-1)$. Biết mặt phẳng $(Q)$ song song với $(P)$ và thỏa mãn khoảng cách từ $A$ đến $(Q)$ bằng khoảng cách giữa hai mặt phẳng $(P)$ và $(Q)$. Hãy xác định cao độ của giao điểm của $(Q)$ với trục $Oz$ (làm tròn kết quả đến hàng phần mười).
\shortans{2,5}
\loigiai{$(Q)\colon x-2y+2z+D=0$ ($D\ne-11$). $\dfrac{|1-2+D|}{3}=\dfrac{|D+11|}{3}\Leftrightarrow D=-5$. $(Q)\colon x-2y+2z-5=0$ cắt $Oz$ tại $z=2{,}5$.}
\end{ex}
